\documentclass[11pt]{article}
\usepackage[margin=1in]{geometry}
\usepackage{amsmath,amssymb,amsthm,listings}
\usepackage[colorlinks=true,urlcolor=blue,linkcolor=blue]{hyperref}
\newtheorem{theorem}{Theorem}
\newtheorem{lemma}{Lemma}
\lstset{language=Python,basicstyle=\ttfamily\small,columns=fullflexible,
 keepspaces=true,showstringspaces=false,breaklines=true}
\title{A counterexample to complete monotonicity\\
in Poisson approximation\\\large MathDB 317893}
\author{Principia Math In Progress}
\date{11 September 2026}
\begin{document}
\maketitle

\section*{Task 1. Resolve the complete-monotonicity conjecture}

For a probability mass function $p$ of mean $\lambda$, write
\[
 D(p)=\sum_{k\geq0}p_k\log
 \frac{p_k}{e^{-\lambda}\lambda^k/k!}.
\]
All logarithms are natural. Yu's Conjecture 1~\cite{Yu} asserts that,
for every fixed $\lambda>0$, both
$D(\operatorname{Bin}(n,\lambda/n))$, $n\geq\lambda$, and
$D(\operatorname{NB}(n,n/(\lambda+n)))$, $n>0$,
are completely monotonic in $n$. For a sequence $a_n$, this means
\[
 (-1)^r\Delta^r a_n
 =\sum_{j=0}^r(-1)^j\binom rj a_{n+j}\geq0
 \quad(r\geq0),\qquad \Delta a_n=a_{n+1}-a_n .
\]

\begin{theorem}
Let $D_n=D(\operatorname{Bin}(n,1/n))$. Then
\begin{equation}
 -\frac{4003}{10^{12}}
 <\sum_{j=0}^{84}(-1)^j\binom{84}{j}D_{4+j}
 <-\frac{4002}{10^{12}}<0.
 \label{eq:counter}
\end{equation}
Consequently the binomial assertion, and therefore the stated conjecture,
is false.
\end{theorem}

We prove the sign by a finite rational certificate. No floating-point
evaluation of the alternating sum is used in the proof.
The negative-binomial assertion is not decided by this counterexample.

\subsection*{An exact expression for the divergence}

Let $K\sim\operatorname{Bin}(n,1/n)$, $n\geq2$, and $Y=n-K$.
Then
\[
 \Pr(Y=y)=\frac{w_{n,y}}{n^n},\qquad
 w_{n,y}=\binom ny(n-1)^y,\qquad
 \sum_{y=0}^n w_{n,y}=n^n .
\]
Dividing the binomial mass at $k$ by the Poisson mass $e^{-1}/k!$
and taking its logarithm gives
\[
 1+\log(n!)-\log((n-k)!)-n\log n+(n-k)\log(n-1).
\]
Since $\mathbb EY=n-1$, averaging yields
\begin{equation}
 D_n=1-n\log n+(n-1)\log(n-1)+\log(n!)
       -\frac1{n^n}\sum_{y=0}^n w_{n,y}\log(y!).
 \label{eq:divergence}
\end{equation}
Here $\log(0!)=\log(1!)=0$.

\subsection*{Rational bounds for all logarithms}

\begin{lemma}
For $1\leq x\leq2$, put $t=(x-1)/(x+1)$ and
\[
 A(x)=2\sum_{k=0}^{59}\frac{t^{2k+1}}{2k+1},
 \qquad
 B(x)=A(x)+\frac{2t^{121}}{121(1-t^2)}.
\]
Then $A(x)\leq\log x\leq B(x)$.
\end{lemma}
\begin{proof}
The identity $\log x=2\operatorname{arctanh}t$ and the geometric
series give
$\log x=2\sum_{k\geq0}t^{2k+1}/(2k+1)$.
For $0\leq t\leq1/3$, the omitted terms are nonnegative and their
sum is at most $2t^{121}/(121(1-t^2))$.
\end{proof}

Set $S=10^{50}$. For $1\leq m\leq88$, write $m=2^e x$ with
$e=\lfloor\log_2m\rfloor$ and $1\leq x<2$, and define the integers
\[
 \ell_m=\left\lfloor S\bigl(A(x)+eA(2)\bigr)\right\rfloor,\qquad
 u_m=\left\lceil S\bigl(B(x)+eB(2)\bigr)\right\rceil.
\]
The exponent $e$ is found by integer comparisons, without evaluating
a logarithm. Thus
\[
 \ell_m\leq S\log m\leq u_m.
\]
Define $F_0=G_0=0$ and
$F_y=\sum_{m=1}^y\ell_m$, $G_y=\sum_{m=1}^y u_m$.
Equation~\eqref{eq:divergence} now gives
\[
 a_n\leq SD_n\leq b_n \qquad(4\leq n\leq88),
\]
where the following floors and ceilings involve rational numbers only:
\begin{align*}
 a_n&=\left\lfloor
 S-nu_n+(n-1)\ell_{n-1}+F_n
 -\frac1{n^n}\sum_{y=0}^n w_{n,y}G_y
 \right\rfloor,\\
 b_n&=\left\lceil
 S-n\ell_n+(n-1)u_{n-1}+G_n
 -\frac1{n^n}\sum_{y=0}^n w_{n,y}F_y
 \right\rceil.
\end{align*}

\subsection*{The integer sign certificate}

Write $c_j=\binom{84}{j}$ and set
\[
 L=\sum_{\substack{0\leq j\leq84\\j\text{ even}}}c_ja_{4+j}
   -\sum_{\substack{0\leq j\leq84\\j\text{ odd}}}c_jb_{4+j},
 \quad
 U=\sum_{\substack{0\leq j\leq84\\j\text{ even}}}c_jb_{4+j}
   -\sum_{\substack{0\leq j\leq84\\j\text{ odd}}}c_ja_{4+j}.
\]
The displayed rational formulas evaluate to
\begin{align*}
 L&=-400267559715299775316466341855642397554631,\\
 U&=-400267559715296274267292737889552448468935.
\end{align*}
Therefore
\[
 -\frac{4003}{10^{12}}
 <\frac LS
 \leq\sum_{j=0}^{84}(-1)^jc_jD_{4+j}
 \leq\frac US
 <-\frac{4002}{10^{12}}.
\]
The inequalities at the two ends are direct integer comparisons.
All indices $4,\ldots,88$ are admissible for the fixed mean $\lambda=1$.
This proves~\eqref{eq:counter}.

For reproducibility, the complete calculation of $L,U$ appears below.
Python's \texttt{Fraction} class performs exact rational arithmetic;
all other operations are on integers. The code implements the formulas
above and requires only the standard library.

\lstinputlisting{certificate_core.py}

\section*{Scope and checks}

The separate verifier \texttt{verify\_counterexample.py} saves every
logarithm interval and divergence interval in \texttt{certificate.json}.
It also recomputes the intervals directly from the original KL sum over
$K$, independently of the expectation formula over $Y$.
The scaled logarithm intervals have width at most $1$; the scaled
divergence intervals have width at most $349$. The large cancellation
in the 84th difference is accounted for by the signed interval sums,
so it cannot turn a numerical rounding error into the certified sign.
No minimality of the order $84$ is asserted.

MathDB and the inspected public claim sources were checked on
11 September 2026. The database remained open with no attached claims,
and no existing matching solution claim was located. The full dated search
record is in \texttt{literature-audit.md}. In particular the later
Lau--Nair--Ng paper~\cite{LNN} explicitly retains the fixed-mean binomial
conjecture; its convexity result changes a different parameter.

\section*{Tasks 2--5}

\begin{enumerate}
\setcounter{enumi}{1}
\item Determine the smallest finite-difference order at which
complete monotonicity fails for some admissible pair $(n,\lambda)$,
and certify an example at that order.
\item Classify the fixed means $\lambda>0$ for which
$D(\operatorname{Bin}(n,\lambda/n))$ fails complete monotonicity,
including whether any positive mean gives a completely monotonic sequence.
\item Determine whether eventual complete monotonicity can hold after
discarding finitely many terms, and relate this to the asymptotics of
high-order finite differences.
\item Prove or refute the remaining negative-binomial assertion
$D(\operatorname{NB}(n,n/(\lambda+n)))$ for fixed $\lambda>0$.
\end{enumerate}

\begin{thebibliography}{9}
\bibitem{Yu}
Y. Yu, \emph{Monotonic convergence in an information-theoretic law of
small numbers}, IEEE Transactions on Information Theory
\textbf{55} (2009), 5412--5422.
\href{https://arxiv.org/abs/0810.5203}{arXiv:0810.5203v4},
Conjecture 1, p.~8.
\bibitem{LNN}
C.~W. Lau, C. Nair and D. Ng,
\emph{A mutual information inequality and some applications},
IEEE Transactions on Information Theory (2023).
\href{https://doi.org/10.1109/TIT.2023.3285928}
{doi:10.1109/TIT.2023.3285928}.
\href{https://davidng-dn.github.io/papers/mutual_information_inequality_journal.pdf}
{Author's journal manuscript}, p.~10.
\end{thebibliography}
\end{document}
